\documentclass[11pt,a4paper]{article}
\usepackage{harmonikabau_preamble}

\newcommand{\hbdoknr}{0007}
\newcommand{\hbtitel}{Treble Reed Block --- Chamber Frequencies D3--C6}
\newcommand{\hbuntertitel}{Seven Variants with Practical Depth Constraint}
\newcommand{\hbheadtext}{Doc.\,0007 --- Treble Reed Block: Chamber Frequencies D3--C6}
\newcommand{\hbfoottext}{7 Variants: A (Ref.) / B,C,Ci (Depth) / D,E,Ei (Opening)}

\AtBeginDocument{%
  \fancyhead[L]{\small\hbheadtext}%
  \fancyhead[R]{\small\thepage}%
  \fancyfoot[C]{\small\color{hb@darkgray}\hbfoottext}%
}

\begin{document}
\hbmaketitle

{\small\itshape Complete analysis for 35~chambers (D3--C6). Seven variants: reference + three curve shapes each (linear, convex, concave) for depth and opening. Accounts for the practical constraint: chamber depth must not fall below the maximum reed amplitude. Amplitude-weighted evaluation ($a_n \sim 1/n$).}

\vspace{2mm}
\begin{center}\small
\begin{tabularx}{\textwidth}{l X}
\hbtoprule
Doc.\,0002 & Flow analysis of bass reed 50\,Hz -- v8 \\
Doc.\,0004 & Frequency variation -- two-filter model \\
Doc.\,0005 & Frequency shift as response indicator \\
Doc.\,0006 & Symbol reference \\
\textbf{Doc.\,0007} & \textbf{Treble reed block -- chamber frequencies (this document)} \\
\midrule
\href{berechnungen_verifikation.py}{berechnungen\_verifikation.py} & Verification script for all calculations \\
\bottomrule
\end{tabularx}
\end{center}
\vspace{4mm}

\section{Geometry and Starting Point}

The treble reed block contains 35~chambers for the notes D3 (146.8\,Hz) to C6 (1046.5\,Hz). Each chamber is a rectangular box with constant width $W = 15{.}5$\,mm and a length that decreases linearly from 50\,mm (D3) to 30\,mm (C6). The outlet opening is circular with $\varnothing$\,10\,mm ($S_\text{neck} = 78{.}5$\,mm$^2$) and a physical neck height of 8\,mm (effective neck length $l_\text{eff} = 16{.}5$\,mm with end correction $0{.}85 \cdot d$).

The Helmholtz frequency of each chamber depends on volume and opening:
\begin{equation}
f_H = \frac{c}{2\pi}\sqrt{\frac{S_\text{neck}}{V \cdot l_\text{eff}}}
\end{equation}

From Doc.\,0005 we know: response problems occur when the 2nd or 3rd overtone of a reed comes close to the chamber resonance $f_H$. This corresponds to $f_H/f < 3$. Above $f_H/f = 4$ the chamber is acoustically ``invisible'' --- coupling through higher overtones ($n \geq 5$) is energetically irrelevant ($a_n^2 < 4\,\%$).

\section{The Practical Constraint: Reed Impact}

The calculation shows that smaller chamber volumes push $f_H$ upward, thereby improving response. But chamber depth cannot be reduced arbitrarily: \textbf{the reed must have room to vibrate.}

A free-beating reed swings down through the slot of the reed plate into the chamber at full volume. If the reed tip touches the chamber floor, the vibration is abruptly stopped --- the reed ``bottoms out''. This produces an ugly clicking sound, damages the reed over time, and limits loudness.

\subsection{Maximum Reed Amplitude in the Treble}

The maximum tip deflection depends on reed length and playing dynamics:

\begin{table}[H]
\centering\small
\resizebox{\linewidth}{!}{\begin{tabular}{l r r r}
\hbtoprule
\textbf{Note range} & \textbf{Reed length} & \textbf{Max.\ amplitude} & \textbf{Min.\ depth} \\
\midrule
D3 (lowest) & $\sim$35\,mm & $\sim$3.5\,mm & 4.0\,mm \\
A3 & $\sim$28\,mm & $\sim$2.8\,mm & 3.3\,mm \\
A4 (middle) & $\sim$18\,mm & $\sim$1.8\,mm & 2.3\,mm \\
A5 & $\sim$10\,mm & $\sim$1.0\,mm & 1.5\,mm \\
C6 (highest) & $\sim$8\,mm & $\sim$0.8\,mm & 1.3\,mm \\
\bottomrule
\end{tabular}}
\caption{Estimated reed amplitudes and minimum depths (amplitude $\approx$ 10\,\% of reed length $+$ 0.5\,mm safety margin)}
\end{table}

Minimum chamber depth must decrease from about 4\,mm (D3) to 1.3\,mm (C6). The variants calculated here with a minimum depth of 3\,mm (C6) are \textbf{well above} the reed-impact limit.

\begin{keybox}
\textbf{Practical rule:} Chamber depth should be as small as possible (minimum volume $\to$ highest $f_H$ $\to$ best decoupling), but \textbf{never} below the maximum reed amplitude plus safety margin. For the lowest treble notes (D3--F3), 3\,mm could be tight --- 4\,mm would be safer.
\end{keybox}

\section{The Three Curve Shapes}

Each geometric parameter (depth or opening diameter) can vary from its starting value $a$ to its final value $b$ in three different ways:
\begin{align}
\text{Linear:}\quad p(x) &= a + x\,(b - a) \\
\text{Convex:}\quad p(x) &= a \cdot (b/a)^x \\
\text{Concave:}\quad p(x) &= a + b - a \cdot (b/a)^{1-x}
\end{align}

All three connect \textbf{the same endpoints}: $p(0) = a$, $p(1) = b$. They differ in the \textbf{distribution} of the reduction across the note range:

\textbf{Linear} distributes the reduction evenly: each semitone step changes the parameter by the same absolute amount (e.g.\ $-0{.}147$\,mm/semitone for $8 \to 3$\,mm).

\textbf{Convex} (lies \textbf{below} the straight line) reduces faster at the beginning and slower at the end. In the middle range the parameter has a \textbf{lower} value than the linear curve. Example depth: at A4, convex gives 4.9\,mm vs.\ linear 5.5\,mm --- 0.6\,mm less.

\textbf{Concave} (lies \textbf{above} the straight line) reduces slower at the beginning and faster at the end. In the middle range the parameter has a \textbf{higher} value than the linear curve. Example depth: at A4, concave gives $\approx$6.1\,mm vs.\ linear 5.5\,mm --- 0.6\,mm more.

\begin{warnbox}
\textbf{Which curve is better depends on which direction the parameter shifts $f_H$:}

If \textbf{less} of the parameter \textbf{raises} $f_H$ (as with depth: less $V \to$ higher $f_H$), the curve that has less in the middle range is better $\to$ \textbf{convex wins}.

If \textbf{more} of the parameter \textbf{raises} $f_H$ (as with opening: more $S \to$ higher $f_H$), the curve that has more in the middle range is better $\to$ \textbf{concave wins}.
\end{warnbox}

\section{The Seven Variants}

\begin{table}[H]
\centering\small
\resizebox{\linewidth}{!}{%
\begin{tabular}{c p{28mm} p{28mm} p{30mm} p{32mm}}
\hbtoprule
\textbf{Var.} & \textbf{Short name} & \textbf{Depth [mm]} & $\boldsymbol{\varnothing}$ \textbf{Opening [mm]} & \textbf{Effect on $f_H$} \\
\midrule
A & Reference & 8 constant & 10 constant & Reference \\
\midrule
B  & d linear & 8$\to$3 linear & 10 constant & $f_H$ rises ($V$ falls) \\
C  & d convex & 8$\to$3 convex & 10 constant & $f_H$ rises more \\
Ci & d concave & 8$\to$3 concave & 10 constant & $f_H$ rises less \\
\midrule
D  & $\varnothing$ linear & 8 constant & 10$\to$6 linear & $f_H$ falls ($S$ falls) \\
E  & $\varnothing$ convex & 8 constant & 10$\to$6 convex & $f_H$ falls more \\
Ei & $\varnothing$ concave & 8 constant & 10$\to$6 concave & $f_H$ falls less \\
\bottomrule
\end{tabular}}
\caption{The seven variants. Upper group: depth varies. Lower group: opening varies.}
\end{table}

\subsection{Why Direction Decides}

\textbf{Reducing depth} decreases chamber volume $V$. Since $f_H \propto 1/\sqrt{V}$, $f_H$ \textbf{rises}. This shifts $f_H/f$ upward --- away from the low overtones. This is the \textbf{right direction}.

\textbf{Reducing opening} decreases neck area $S_\text{neck}$ and also the end correction ($l_\text{eff} = 8 + 0{.}85 d$ also falls). Net effect: $f_H \propto \sqrt{S/l_\text{eff}}$ \textbf{falls}. Quantitatively: $\varnothing$ from 10 to 6\,mm lowers $f_H$ to \textbf{67\,\%}. This shifts $f_H/f$ \textbf{downward} --- the \textbf{wrong direction}.

\section{Main Result}

\begin{table}[H]
\centering\small
\resizebox{\linewidth}{!}{\begin{tabular}{l r r r r r r r}
\hbtoprule
& \textbf{A} & \textbf{B} & \textbf{C} & \textbf{Ci} & \textbf{D} & \textbf{E} & \textbf{Ei} \\
\midrule
$f_H$ at D3 [Hz] & 1513 & 1513 & 1513 & 1513 & 1513 & 1513 & 1513 \\
$f_H$ at C6 [Hz] & 1953 & 3189 & 3189 & 3189 & 1315 & 1315 & 1315 \\
$f_H/f$ min. & 1.9 & 3.0 & 3.0 & 3.0 & 1.3 & 1.3 & 1.3 \\
\midrule
Critical (OT $\leq 2$) & 7 & 0 & \textbf{0} & 0 & 12 & 13 & 12 \\
Noticeable (OT $= 3$) & 7 & 7 & \textbf{5} & 8 & 5 & 5 & 5 \\
Good (OT $\geq 4$) & 21 & 28 & \textbf{30} & 27 & 18 & 17 & 18 \\
$\sum R_\text{eff}$ & 17.4 & 12.8 & \textbf{12.3} & 13.5 & 17.2 & 17.3 & 17.0 \\
\bottomrule
\end{tabular}}
\caption{Main comparison. Bold: variant C (best result).}
\end{table}

\subsection{Ranking}

\begin{enumerate}
\item \textbf{C} (depth convex): 0~crit., 5~noticeable, 30~good. $\sum R = 12{.}3$
\item \textbf{B} (depth linear): 0~crit., 7~noticeable, 28~good. $\sum R = 12{.}8$
\item \textbf{Ci} (depth concave): 0~crit., 8~noticeable, 27~good. $\sum R = 13{.}5$
\item \textbf{Ei} ($\varnothing$ concave): 12~crit., 5~noticeable, 18~good. $\sum R = 17{.}0$
\item \textbf{D} ($\varnothing$ linear): 12~crit., 5~noticeable, 18~good. $\sum R = 17{.}2$
\item \textbf{E} ($\varnothing$ convex): 13~crit., 5~noticeable, 17~good. $\sum R = 17{.}3$
\item \textbf{A} (reference): 7~crit., 7~noticeable, 21~good. $\sum R = 17{.}4$
\end{enumerate}

\section{Analysis of the Depth Group}

All three depth variants (B, C, Ci) eliminate critical resonances completely (0~critical notes). They differ only in the number of noticeable resonances:

\textbf{C (convex, below the straight line):} 5~noticeable. The curve gives the middle notes less volume $\to$ higher $f_H$ $\to$ better decoupling. Example A4: depth 4.9\,mm, $f_H/f = 5{.}1$. \textbf{Best depth variant.}

\textbf{B (linear):} 7~noticeable. Uniform reduction. A4: depth 5.5\,mm, $f_H/f = 4{.}8$.

\textbf{Ci (concave, above the straight line):} 8~noticeable. The curve retains more volume in the middle range $\to$ lower $f_H$. Worst depth variant --- but still 0~critical.

The difference between the three depth variants is \textbf{moderate} (12.3 vs.\ 12.8 vs.\ 13.5 in $\sum R_\text{eff}$). The \textbf{big jump} is from reference A (17.4) to any of the three depth variants --- depth reduction itself is much more important than the curve shape.

\begin{keybox}
\textbf{Practical significance:} The difference between convex and linear is 2~fewer noticeable resonances --- audible, but not dramatic. The difference between depth reduction and none (A vs.\ B/C/Ci) is 7~critical resonances --- this is the decisive step. \textbf{Reducing the depth achieves 95\,\% of the benefit. Curve shape is fine-tuning.}
\end{keybox}

\section{Analysis of the Opening Group}

The three opening variants (D, E, Ei) make the situation worse compared to the reference --- they have \textbf{more} critical notes (12--13 vs.\ 7). The reason: $\varnothing$ from 10 to 6\,mm reduces $S_\text{neck}$ from 78.5 to 28.3\,mm$^2$ ($-$64\,\%) and $l_\text{eff}$ from 16.5 to 13.1\,mm ($-$21\,\%). Net effect: $f_H$ falls to \textbf{67\,\%}. $f_H/f$ drops to 1.3 at C6 --- there $f_H$ lies \textbf{below} the reed frequency itself.

\textbf{Ei (concave):} 12~crit. --- least harmful, because the curve retains the larger openings in the middle range.

\textbf{E (convex):} 13~crit. --- most harmful, because the curve reduces the opening most in the middle range.

\begin{warnbox}
\textbf{Conclusion, opening group:} None of the three opening variants is better than reference A. Reducing the opening from 10 to 6\,mm alone is \textbf{counterproductive}. Even the best opening variant (Ei) has 12~critical notes --- nearly twice as many as the reference (7). Varying the opening alone is the wrong approach.
\end{warnbox}

\section{Minimise Volume --- The Practical Optimisation Rule}

The analysis yields a clear optimisation rule, consistent with the practical experience of instrument builders:

\begin{keybox}
\textbf{Optimisation rule: volume as small as the reed allows.}

For each chamber, depth is chosen such that:
\begin{enumerate}
\item The reed does \textbf{not bottom out} at maximum volume (depth $\geq$ max.\ amplitude $+$ 0.5\,mm safety margin).
\item Volume is \textbf{no larger} than required by condition (1).
\item The opening remains \textbf{as large as possible} (10\,mm), because any reduction pushes $f_H$ in the wrong direction.
\end{enumerate}
\end{keybox}

In practice: chamber depth follows reed amplitude --- short reeds (high notes) need less space, receive shallower chambers and automatically benefit from higher $f_H/f$. The convex depth profile (C) is slightly better than linear (B) because it reduces more aggressively in the middle range.

The result agrees with instrument-building practice: experienced builders make chambers just deep enough for the reed not to bottom out, and no deeper. The calculation now shows \textbf{why} this works: less volume $\to$ higher $f_H$ $\to$ fewer overtone resonances $\to$ better response.

\section{Summary}

\begin{keybox}
\textbf{Seven variants for the treble reed block D3--C6:}

\textbf{Reducing depth} (B, C, Ci) raises $f_H$ --- right direction, 0~critical notes.\\
\textbf{Reducing opening} (D, E, Ei) lowers $f_H$ --- wrong direction, 12--13~critical.

Within depth: \textbf{convex $>$ linear $>$ concave} (12.3 vs.\ 12.8 vs.\ 13.5).\\
Within opening: \textbf{concave $>$ linear $>$ convex} (17.0 vs.\ 17.2 vs.\ 17.3).

In both cases, the curve wins that retains more of the parameter that pushes $f_H$ upward.

\textbf{The big jump} lies in depth reduction itself (7~fewer critical), not in the curve shape (2~fewer noticeable). Reducing depth achieves 95\,\% of the benefit.

\textbf{Practical rule:} Chamber depth $=$ maximum reed amplitude $+$ 0.5\,mm. No deeper. Opening as large as possible.

Keep $f_H/f \geq 3$ $\to$ no 2nd OT hit $\to$ no critical problems.
\end{keybox}

\vspace{3mm}
\begin{center}
\textcolor{accentred}{\rule{0.6\textwidth}{1pt}}\\[4pt]
{\small\itshape Minimum volume, maximum opening --- the chamber as small as the reed allows.}
\end{center}


\end{document}
